% Formules des courbes d'analyse des placements financiers (Budget App).
% Compilation : pdflatex formules-courbes-analyse-placements.tex   (deux passes pour la table des matières)
%           ou : xelatex  formules-courbes-analyse-placements.tex
\documentclass[11pt,a4paper]{article}

\usepackage[utf8]{inputenc}
\usepackage[T1]{fontenc}
\usepackage{lmodern}
\usepackage[french]{babel}
\usepackage{amsmath,amssymb}
\usepackage{booktabs}
\usepackage{array}
\usepackage{xcolor}
\usepackage{enumitem}
\usepackage[a4paper,margin=2.3cm]{geometry}
\usepackage{microtype}
\usepackage[hidelinks]{hyperref}

\definecolor{bleu}{HTML}{1D4ED8}
\definecolor{fond}{HTML}{EFF6FF}

\usepackage{titlesec}
\titleformat{\section}{\Large\bfseries\color{bleu}}{\thesection}{0.7em}{}
\titleformat{\subsection}{\large\bfseries\color{bleu}}{\thesubsection}{0.6em}{}

% Un encadré de formules.
\usepackage[most]{tcolorbox}
\newtcolorbox{formules}{colback=fond,colframe=bleu,boxrule=0.6pt,arc=2pt,
  left=6pt,right=6pt,top=4pt,bottom=4pt,before skip=8pt,after skip=8pt}

\setlist[itemize]{itemsep=2pt,topsep=3pt}

\newcommand{\avant}{\mathrm{avant}}

\title{\textbf{\color{bleu}Courbes d'analyse des placements financiers}\\[4pt]
  \large Formules de calcul exactes des deux mesures du menu \og Courbes par élément \fg}
\date{}
\author{}

\begin{document}
\maketitle

\noindent\textit{Sources : \texttt{backend/app/services/evolution\_portefeuille.py} (calcul serveur) et
\texttt{extensions/placements/evolution.js} (rebasage sur la période affichée).}

\medskip
\noindent L'onglet Analyse trace d'abord la \textbf{valeur brute du portefeuille} (courbe du total), puis, quand un
\og Séparer par \fg{} est actif, \textbf{un graphe par élément} selon l'une des deux mesures ci-dessous.

\tableofcontents

% ---------------------------------------------------------------------------
\section{Éléments communs aux deux mesures}

\begin{center}
\small
\begin{tabular}{@{}>{\raggedright\arraybackslash}p{2.6cm}>{\raggedright\arraybackslash}p{12.6cm}@{}}
\toprule
\textbf{Notation} & \textbf{Définition}\\
\midrule
$P(j)$ & Dernier cours connu du titre à la date $j$ ou avant. Sources, de la plus à la moins fiable : l'historique des
 cours (table \texttt{cours\_titre}), puis les prix des mouvements (achats et ventes), puis le cours actuel pour
 aujourd'hui. Avant son premier cours connu, un titre n'a pas de valeur.\\[3pt]
$Q(j)$ & Quantité détenue à la date $j$, déroulée depuis les mouvements (achat $+$, vente $-$), jamais négative.\\[3pt]
$\mathrm{part}$ & Poids du titre dans l'élément tracé. Vaut $1$ pour les axes total, compte et titre ; pour
 enveloppe, classe, zone et secteur, c'est le poids d'exposition du titre (un ETF \og Monde \fg{} se répartit entre
 plusieurs zones), \textbf{sans renormalisation}.\\[3pt]
$V(j)$ & Valeur de l'élément à la date $j$.\\[3pt]
Grille de dates & Tous les jours tant qu'il y a au plus $200$ points, sinon un pas régulier ; le dernier point est
 toujours aujourd'hui. Une \og étape \fg{} est un pas de cette grille, pas forcément un jour.\\[3pt]
Monnaie & Un graphe par monnaie : deux monnaies ne s'additionnent jamais.\\
\bottomrule
\end{tabular}
\end{center}

\begin{formules}
\[
V(j)\;=\;\sum_{\text{titres}} Q(j)\times P(j)\times \mathrm{part}
\]
\end{formules}

% ---------------------------------------------------------------------------
\section{Mesure \og Valeur$(t)$ / montant investi$(t)$ \fg{} (choix par défaut)}

C'est la \textbf{plus-value latente en pourcentage}, celle des indicateurs de l'onglet Portefeuille. Ce n'est pas un
rendement pondéré par le temps.

\begin{formules}
\[
\boxed{\;\mathrm{rendement}(t)\;=\;\frac{V(t)}{I(t)}\;-\;1\;}
\]
\end{formules}

$I(t)$ est le \textbf{montant investi} : le coût des titres encore détenus, au prix de revient moyen pondéré, tenu par
couple (compte, titre). Il évolue ainsi à chaque mouvement :

\begin{formules}
\begin{align*}
\text{Achat :}\quad & I \;\leftarrow\; I + q_{\text{achetée}}\times \mathrm{prix}_{\text{unitaire}}\\[2pt]
\text{Vente :}\quad & I \;\leftarrow\; I\times\Bigl(1-\min\bigl(1\,;\,q_{\text{vendue}}\,/\,q_{\text{détenue avant la vente}}\bigr)\Bigr)
\end{align*}
\end{formules}

\begin{itemize}
  \item Quand la quantité détenue devient nulle, $I$ est remis à $0$.
  \item Si $I(t)\le 10^{-9}$ (rien d'investi à cette date), il n'y a pas de dénominateur : la valeur est \og vide \fg{}
        et la courbe est \textbf{interrompue}.
  \item Le choix de la période (1 mois, 6 mois, 1 an, Tout) découpe simplement la courbe : \textbf{aucun rebasage},
        chaque point reste $V/I-1$ à sa propre date.
  \item Le texte sous la carte affiche la plus-value latente $V-I$ (montant), puis le pourcentage, puis
        \og investi $I$ \fg.
\end{itemize}

% ---------------------------------------------------------------------------
\section{Mesure \og Cours \fg{} (performance des cours seuls)}

Elle ne retient que l'effet de la variation des cours : \textbf{l'argent ajouté ou retiré par un achat ou une vente
n'y entre jamais}. C'est un rendement pondéré par le temps (\emph{time-weighted}).

\subsection*{a) Gain de chaque étape $k$}

Avec $j_k$ la date de l'étape, $j_{k-1}$ la précédente et $Q_{\avant}$ la quantité détenue \emph{avant} les
mouvements de l'étape :

\begin{formules}
\[
\boxed{\;g_k\;=\;\sum Q_{\avant}\times\bigl(P(j_k)-P(j_{k-1})\bigr)\times \mathrm{part}\;}
\]
\end{formules}

Un achat du jour n'a donc rien gagné, et une vente du jour a gagné jusqu'à la veille.

\subsection*{b) Gain cumulé et rendement chaîné}

Avec $V_{i-1}$ la valeur de la série à l'étape précédente :

\begin{formules}
\[
G_k\;=\;\sum_{i\le k} g_i
\qquad\qquad
R_k\;=\;\prod_{i\le k}\Bigl(1+\frac{g_i}{V_{i-1}}\Bigr)-1
\]
\end{formules}

\begin{itemize}
  \item Si $V_{i-1}\le 10^{-9}$ (la série n'existait pas la veille), le facteur de l'étape vaut $1$, soit $0\,\%$.
        Au tout premier point la base est $0$ : $R_0=0$.
  \item Pour le \textbf{total} du portefeuille, le gain d'une étape est la somme des gains des séries, et le
        rendement se chaîne avec la valeur totale de la veille comme base.
\end{itemize}

\subsection*{c) Courbe affichée : rebasée au début de la fenêtre}

Avec $d$ le premier point visible :

\begin{formules}
\[
\boxed{\;\mathrm{courbe}(k)\;=\;\frac{1+R_k}{1+R_d}-1\;}
\]
\end{formules}

\subsection*{d) Texte \og Performance des cours \fg{} sur la fenêtre}

\begin{formules}
\[
\mathrm{gain}=G_{\mathrm{fin}}-G_d
\qquad\qquad
\mathrm{pourcentage}=\frac{1+R_{\mathrm{fin}}}{1+R_d}-1
\]
\end{formules}

% ---------------------------------------------------------------------------
\section{Comment choisir entre les deux}

\begin{center}
\small
\begin{tabular}{@{}>{\raggedright\arraybackslash}p{3.1cm}>{\raggedright\arraybackslash}p{6.0cm}>{\raggedright\arraybackslash}p{6.0cm}@{}}
\toprule
 & \textbf{Valeur$(t)$ / investi$(t)$} & \textbf{Cours (performance)}\\
\midrule
Question posée & \og Où en est ma plus-value sur ce que j'ai mis dans cet élément ? \fg
 & \og Les cours de mes titres ont-ils monté ou baissé, indépendamment de mes achats et ventes ? \fg\\[3pt]
Sensible aux achats / ventes & Oui : le prix d'achat fixe le coût $I$. & Non : seuls les cours comptent.\\[3pt]
Rebasage sur la période & Non. & Oui : $0\,\%$ au premier point visible.\\[3pt]
Calcul & Rapport de deux niveaux à la même date. & Produit des rendements d'étape (chaîné).\\[3pt]
Trou dans la courbe & Quand rien n'est investi. & Jamais ; une étape sans base compte $0\,\%$.\\
\bottomrule
\end{tabular}
\end{center}

% ---------------------------------------------------------------------------
\section{Exemple chiffré (un seul titre, axe \og Titre individuel \fg)}

Scénario : on achète $10$ titres à $100\,€$, le cours monte à $110\,€$, on en achète $10$ de plus à $110\,€$, puis le
cours monte de nouveau de $10\,\%$ ($121\,€$).

\begin{center}
\small
\begin{tabular}{@{}cllrrr@{}}
\toprule
\textbf{Étape} & \textbf{Événement} & \textbf{Cours} & $Q$ & $V$ & $I$\\
\midrule
0 & Achat de 10 à 100\,€ & 100 & 10 & 1\,000 & 1\,000\\
1 & Le cours monte & 110 & 10 & 1\,100 & 1\,000\\
2 & Achat de 10 à 110\,€ & 110 & 20 & 2\,200 & 2\,100\\
3 & Le cours monte de 10\,\% & 121 & 20 & 2\,420 & 2\,100\\
\bottomrule
\end{tabular}
\end{center}

\subsection*{Mesure \og Valeur / investi \fg{} à l'étape 3}
\begin{formules}
\[
\frac{2\,420}{2\,100}-1 \;=\; \mathbf{+15{,}24\,\%}\qquad(\text{plus-value latente : }320\,€)
\]
\end{formules}

\subsection*{Mesure \og Cours \fg{} à l'étape 3}
\begin{center}
\small
\begin{tabular}{@{}cccccrr@{}}
\toprule
\textbf{Étape} & $Q_{\avant}$ & $g_k$ & Base $V_{k-1}$ & Facteur & $G_k$ & $R_k$\\
\midrule
0 & --- & 0 & 0 & 1 & 0 & $0\,\%$\\
1 & 10 & $10\times(110-100)=100$ & 1\,000 & $1+100/1\,000=1{,}10$ & 100 & $+10\,\%$\\
2 & 10 & $10\times(110-110)=0$ & 1\,100 & $1{,}00$ & 100 & $+10\,\%$\\
3 & 20 & $20\times(121-110)=220$ & 2\,200 & $1+220/2\,200=1{,}10$ & 320 & $+21\,\%$\\
\bottomrule
\end{tabular}
\end{center}

\begin{formules}
\[
R_3 \;=\; 1{,}10\times 1{,}00\times 1{,}10-1 \;=\; \mathbf{+21\,\%}\qquad(G_3=320\,€)
\]
\end{formules}

\begin{itemize}
  \item Les deux mesures diffèrent ($15{,}24\,\%$ contre $21\,\%$) parce que le second achat a été fait à un cours
        déjà monté : il alourdit le montant investi, mais pas la performance des cours.
  \item \textbf{Rebasage.} Si la fenêtre affichée commence à l'étape 2 ($R_d=+10\,\%$), la courbe \og Cours \fg{}
        vaut $0\,\%$ à l'étape 2, puis $1{,}21/1{,}10-1=\mathbf{+10\,\%}$ à l'étape 3. La mesure \og Valeur / investi
        \fg{} reste à $+15{,}24\,\%$ : elle ne se rebase pas.
  \item \textbf{Vente.} Si l'on vendait ensuite $5$ titres à $120\,€$ : le coût passe de $2\,100$ à
        $2\,100\times(1-5/20)=1\,575\,€$, la quantité tombe à $15$, $V=15\times120=1\,800\,€$ et
        \og Valeur / investi \fg{} $=1\,800/1\,575-1=\mathbf{+14{,}29\,\%}$.
\end{itemize}

% ---------------------------------------------------------------------------
\section{Les barres rouges : achats et ventes}

Sur les graphes d'analyse, une \textbf{barre verticale rouge en pointillés} marque chaque étape de la grille où au
moins un achat ou une vente a eu lieu. Son infobulle donne le montant acheté et vendu ; un bouton ouvre la liste de
chaque ligne responsable.

L'effet d'un mouvement se mesure en comparant l'\emph{après} à l'\emph{avant}, \textbf{au cours du jour} (la variation
du cours n'est pas le fait du mouvement). Avec $V^{-}$ et $I^{-}$ la valeur et le montant investi juste avant les
mouvements de l'étape, $V^{+}$ et $I^{+}$ juste après :

\begin{formules}
\[
\text{effet sur la valeur} = V^{+}-V^{-}
\qquad
\text{en \%} = \frac{V^{+}-V^{-}}{V^{-}}
\]
\[
\text{variation due au mouvement (en points)} \;=\; \Bigl(\frac{V^{+}}{I^{+}}-1\Bigr)-\Bigl(\frac{V^{-}}{I^{-}}-1\Bigr)
\]
\end{formules}

\begin{itemize}
  \item Courbe en montant (total) : l'infobulle donne l'effet sur la valeur, en montant et en \%.
  \item Mesure \og Valeur / investi \fg{} : elle donne la variation due au mouvement, en points de pourcentage
        (absente pour une première position : il n'y a pas de pourcentage avant).
  \item Mesure \og Cours \fg{} : le mouvement est \textbf{sans effet} sur la performance des cours, par construction.
  \item Une ligne qui n'appartient qu'en partie à l'élément (ETF réparti sur des zones ou des secteurs) affiche son
        montant total, sa part et le \textbf{montant alloué} ($\text{montant}\times\mathrm{part}$).
\end{itemize}

% ---------------------------------------------------------------------------
\section{Cas particuliers}

\begin{itemize}
  \item \textbf{Sans historique de cours} : un élément qui a moins de deux cours rangés en base n'a pas de courbe,
        seulement la carte \og Pas d'historique de valeurs disponible \fg. Les prix des mouvements servent à tracer
        mais ne comptent pas comme historique.
  \item \textbf{Titre individuel} : seuls les titres cochés sont tracés ; mêmes formules, titre par titre.
  \item \textbf{Axes enveloppe, classe, zone, secteur} : un titre est réparti selon ses poids d'exposition
        ($\mathrm{part}$) ; ce qui n'est pas renseigné tombe dans \og Non renseigné \fg, sans renormaliser.
\end{itemize}

\end{document}
